\section{Related work}
\paragraph{Low-latency enhancement.}
Conv-TasNet~\cite{luo2019convtasnet} replaced the STFT with a short learned
filterbank, which made millisecond-scale latency possible in a time-domain
separator. Its causal variant reports a single operating point (2\,ms at
8\,kHz). STFT-domain systems reach low latency with asymmetric analysis and
synthesis windows~\cite{mauler2007lowdelay}. Wang et al.~\cite{wang2023lowlatency}
showed that a DNN with a long asymmetric analysis window and a 2--4\,ms
synthesis window keeps most of the quality of a 32\,ms system.
DeepFilterNet~\cite{schroeter2022deepfilternet} adds complex filtering over
neighbouring frames, which we restrict to \emph{past} frames.

\paragraph{Binaural enhancement and cue preservation.}
Classical binaural noise reduction keeps cues by applying one real gain to
both ears, or by constraining the binaural multichannel Wiener
filter~\cite{doclo2015assisted}. Han et al.~\cite{han2020binaural}
preserved spatial cues in a real-time binaural Conv-TasNet, and Tokala et
al.~\cite{tokala2024binaural} trained with explicit ILD/IPD losses. Both
report one latency.

\paragraph{What is new here.}
A single controlled sweep that separates three costs:
\begin{itemize}
\item the window (an oracle bound),
\item the future (causal vs.\ non-causal twins), and
\item learning (oracle vs.\ non-causal model).
\end{itemize}
The sweep also includes a fixed-budget window-vs-lookahead split, target
and residual-noise cue errors against JNDs, and measured single-thread
real-time factors with streaming outputs verified to equal the offline
ones.
